% Generated draft: requires mathematical and source review.
\documentclass[11pt]{article}
\usepackage[margin=25mm]{geometry}
\usepackage{fontspec}
\setmainfont{Latin Modern Roman}
\usepackage{amsmath,amssymb,mathtools,unicode-math}
\setmathfont{Latin Modern Math}
\usepackage{enumitem,needspace,xurl,hyperref}
\setlength{\parindent}{0pt}
\setlength{\parskip}{5pt}
\setlength{\emergencystretch}{3em}
\newcommand{\N}{\mathbb{N}}
\newcommand{\Z}{\mathbb{Z}}
\newcommand{\Q}{\mathbb{Q}}
\newcommand{\R}{\mathbb{R}}
\newcommand{\C}{\mathbb{C}}
\newcommand{\F}{\mathbb{F}}
\newcommand{\A}{\mathbb{A}}
\newcommand{\PP}{\mathbb{P}}
\newcommand{\E}{\mathbb{E}}
\newcommand{\eps}{\varepsilon}
\newcommand{\dd}{\,\mathrm{d}}
\begin{document}
\section*{UnsolvedMath ID: 30006357}
\textbf{Source problem code:} catalogue entry
\subsection{Definitions and mathematical statement}
Under the customary interpretation in terms of finite-index subfactors of type \(\mathrm{II}_1\) factors, let \(N\subset M\) be a unital inclusion and let \([M:N]=\dim_N L^2(M)\) denote its Jones index. The inclusion is irreducible if \(N'\cap M=\C 1\). A factor is hyperfinite if it is the weak operator closure of a directed increasing union of finite-dimensional unital \(*\)-subalgebras. Finite depth means that only finitely many isomorphism classes of irreducible bimodules occur in the standard invariant generated by the inclusion, equivalently that its principal graph is finite.

For the purpose of stating a candidate interpretation, let an \(A_\infty\)-subfactor mean a finite-index inclusion \(P\subset Q\) of type \(\mathrm{II}_1\) factors whose principal graph is the rooted bipartite graph
\[
V=\{0,1,2,\ldots\},\qquad E=\bigl\{\{n,n+1\}:n\geq 0\bigr\},
\]
rooted at \(0\), with bipartition given by parity. The supplied record does not define this terminology, so this interpretation requires confirmation.

\textbf{Supplied Conjecture 1, under this interpretation.} For every finite-index inclusion \(N\subset M\) of hyperfinite type \(\mathrm{II}_1\) factors, if \(N'\cap M=\C1\) and \(N\subset M\) has finite depth, then there exist type \(\mathrm{II}_1\) factors \(P\subset Q\) such that
\[
\Gamma(P\subset Q)\cong A_\infty,
\qquad [Q:P]=[M:N],
\]
where \(\Gamma\) denotes the rooted principal graph.

\textbf{Unresolved scope qualifications.} The source does not explicitly specify the factor type, separability conventions, or whether the realizing inclusion \(P\subset Q\) must itself be hyperfinite. The displayed formulation leaves hyperfiniteness of the realizing inclusion unresolved. More substantially, the supplied universal assertion includes no lower bound on \([M:N]\). Under the graph-based meaning of \(A_\infty\), compatibility with low-index cases requires clarification; a restriction such as \([M:N]>4\) must not be inserted without source confirmation. Consequently, the displayed assertion is a transcription of the supplied target under stated conventions, not a mathematically certified formulation.

\textbf{Supplied partial-results context.} The record labels the problem partially solved and reports that several exceptional indices have been treated. Its inline list has corrupted radicals and fractions. The accompanying table, captioned as concerning hyperfinite finite-depth subfactors of index less than \(5.25\), appears to read as follows; these are provisional readings of the supplied extraction:
\[
\begin{array}{c|c}
\text{Index}&\text{Reported number of subfactors}\\\hline
(5+\sqrt{13})/2&2\\
\approx 4.3772&2\\
(5+\sqrt{17})/2&2\\
3+\sqrt{3}&2\\
(5+\sqrt{21})/2&2\\
5&7\\
\approx 5.04892&2\\
3+\sqrt{5}&11
\end{array}
\]
The table's counting conventions and completeness qualifications are not supplied. The inline assertion of established cases appears to include these displayed indices except the row \(\approx5.04892\); table membership alone must not be read as an assertion that the conjecture was established at that index. No exact algebraic values are supplied for the decimal entries. The source describes combining a classical construction of hyperfinite subfactors with planar algebras and fusion-category techniques.
\subsection{Short English statement}
The supplied conjecture asks whether every index of an irreducible hyperfinite finite-depth subfactor can also occur for a subfactor with principal graph \(A_\infty\). Its intended index range and the hyperfiniteness requirement for the realizing subfactor need clarification. The supplied partially solved status and exceptional-index claims are dated curation, not independently verified results.
\subsection{Sources}
\begin{itemize}
\item Local registry: \texttt{data/unsolvedmath/problems.json}, with \texttt{catalog/problems.json} as the additional-record registry; supplied permanent UnsolvedMath ID \(30006357\), exact source problem code \texttt{catalogue entry}, title ``A-Infinity Realization of Finite-Depth Subfactor Indices.'' These registry identifiers are supplied provenance; no local-file inspection is claimed.
\item Supplied primary report citation: \emph{Subfactors and Applications (2026)}, Oberwolfach Reports, DOI \texttt{10.4171/owr/2025/34}, pp.\ 1823--1886. \url{https://doi.org/10.4171/owr/2025/34}. The extracted assertion is labelled ``Conjecture 1.'' No individual proposer is identified; the supplied proposed year is 2025. The title year 2026 and the DOI/report year 2025 are retained without reconciliation.
\item Supplied literature source: Bisch--Caceres, \emph{New hyperfinite subfactors with infinite depth}. \url{https://arxiv.org/abs/2509.00227}.
\item Supplied literature source: \emph{Oberwolfach Report 34/2025}. \url{https://ems.press/content/serial-article-files/52239}.
\item Literature assessment dated exactly \texttt{2026-08-22}: the supplied curation labels the problem partially solved, reports verification at several exceptional indices, and states that the universal assertion remains unproved. This is supplied dated curation derived from the Oberwolfach material, not an independently verified literature assessment. The record explicitly calls for independent verification of time-sensitive status and cited claims.
\item The registry labels the statement \texttt{corrected\_verified} and says that only formula delimiters, notation, grammar, and typography were repaired while preserving mathematical scope. This supplied editorial label does not resolve the formulation issues identified here.
\end{itemize}
\end{document}
